\documentclass[aps,prx,reprint,superscriptaddress,nofootinbib,floatfix]{revtex4-2}

\usepackage{graphicx}
\usepackage{amsmath,amssymb}
\usepackage{bm}
\usepackage{xcolor}
\usepackage[colorlinks=true,allcolors=blue]{hyperref}
\usepackage{mhchem}
\usepackage{float}

\graphicspath{{figures/}{../figures/}}

\begin{document}

\title{What a radical-pair quantum reservoir would require:\\
readable capacity, no quantum advantage, and a nuclear-register criterion}

\author{Hikaru Wakaura}
\author{Taiki Tanimae}
\affiliation{QIRI (Quantum Integrated Research Institute Inc.),
1--16--3, Akasaka, Minato-ku, Tokyo 107--0061, Japan.\\
Correspondence:\ h.wakaura@qiri.co.jp
(ORCID 0000-0001-8381-8323), t.tanimae@qiri.co.jp}  
  
 
\date{\today}     

\begin{abstract} 
The radical-pair mechanism, by which a field as weak as the Earth's alters a
chemical reaction, underlies the leading model of avian magnetoreception. Whether
the electron--nuclear spin coherence behind this sensing could also support
biological \emph{computation}---a ``quantum brain''---has been conjectured for
decades but never made quantitative. We treat a cryptochrome-type radical pair as a
reservoir computer, solving the radical-pair master equation with electronic
decoherence and driving it as a real turnover does---by the birth of a
spin-correlated pair whose singlet character encodes the input. Three results
follow, none of them the expected one. First, the reservoir is \emph{readable by
ordinary chemistry}: the time-resolved yield of one product pool gives an
out-of-sample information-processing capacity of $2.0$, and the nuclear
polarisation the product carries away raises it to $4.7$, with no spectroscopy
whatever. Second, there is no quantum advantage: a
classical network given the same number of readout features matches the five-spin
reservoir. Third, and decisively, the memory horizon is set \emph{not} by the
microsecond radical-pair lifetime but by the turnover interval, because the nuclear
register survives recombination in the diamagnetic product. A physiological
$1$--$100$~ms turnover changes the memory capacity by $<0.2\%$ and places the
horizon at $3$--$280$~ms, inside the neural band. The binding constraint is
therefore chemical, not temporal, and falsifiable: the \emph{same} nuclear register
must be reused from one turnover to the next, failing which the capacity collapses
to a memoryless read-back. This replaces a timescale argument, which needs no
quantum mechanics, with a criterion that does.
\end{abstract}  

\maketitle 

\section{Introduction}  

The radical-pair mechanism (RPM) is the only experimentally established way in
which a magnetic field as weak as the geomagnetic field ($\sim50~\mu$T) can change
the outcome of a chemical reaction~\cite{Steiner1989,Hore2016,Timmel1998}. A spin-correlated
pair born in a singlet electronic state undergoes coherent singlet--triplet (S--T)
interconversion driven by hyperfine fields; because recombination is
spin-selective, a field that re-tunes this interconversion alters the product
yields. The mechanism underlies the leading model of avian magnetoreception, where
the photoreceptor cryptochrome hosts a flavin--tryptophan radical
pair~\cite{Ritz2000,Maeda2012,Hore2016,Xu2021,Hiscock2016}.

The success of the RPM has repeatedly motivated a bolder conjecture: that
spin coherence in biomolecules might support not merely sensing but
\emph{information processing}---a ``quantum brain''~\cite{Fisher2015,Kerskens2022},
a proposal long contested on decoherence grounds~\cite{Tegmark2000} and part of the
broader programme of quantum biology~\cite{Lambert2013,Cao2020}.
A concrete, named instance is the \emph{three-layer quantum-brain
hypothesis}~\cite{Wakaura2025,Wakaura2026}, which posits a layered architecture---a
long-lived nuclear-spin memory (Layer~1), a radical-pair read/write interface
(Layer~2), and classical neuronal transduction (Layer~3)---and which, in its most
recent form, already reframes the proposal from quantum \emph{computation} to
quantum \emph{reservoir dynamics}. The present paper is the quantitative test of
that architecture as a reservoir, and the bound it returns. A companion study from
this program places the conjecture on quantitative footing from the sensing side. Solving the radical-pair master equation with electronic
decoherence explicitly included, it finds that the active-site radical pairs of
the principal neurotransmitter-metabolising flavoenzymes (monoamine oxidase A/B,
D-amino-acid oxidase) are magnetically inert ($<10^{-4}\%$ field effect): their
large exchange coupling ($|J|\sim0.75$~GHz) and picosecond lifetimes each
independently forbid weak-field S--T mixing. Only a \emph{spatially separated}
radical pair, as in cryptochrome---with $J\approx0$ and a microsecond
lifetime---sustains the coherence required for weak-field
sensitivity~\cite{Efimova2008,Companion2026}. This premise is a standard result of
radical-pair compass theory~\cite{Efimova2008,Hore2016} and is logically
independent of the companion analysis; we take the existence of a long-coherence
separated radical pair as our self-contained starting point.

This sharp dichotomy invites the question we address here. The separated radical
pair is singled out because it preserves long-lived electron--nuclear spin
coherence. \emph{Is that coherence good for anything besides sensing the
geomagnetic field?} We argue the question is best posed not as gate-model quantum
computation but as \emph{reservoir} computing, and we answer it with the chemistry
left in: driven, read and clocked as a turnover actually is. We state at the outset
the load-bearing caveat: whether a \emph{neuronal} cryptochrome forms a
light- or reaction-driven separated radical pair at all is unestablished---the
mammalian CRY1/CRY2 proteins are primarily circadian transcriptional
regulators~\cite{Partch2014}---so every biological conclusion below is conditional
on the existence of such a pair. The spin-dynamics results (capacity, classical
baselines, the register criterion) do not depend on that condition and stand on
their own.

\paragraph*{Why a reservoir, not a gate.}
Gate-model quantum computation treats decoherence as the adversary and demands
error correction; a warm, wet, $10^{10}$-copy biomolecular ensemble cannot meet
that bar. Reservoir computing (RC) inverts the relationship. An RC system feeds a
time-dependent input into a fixed nonlinear dynamical ``reservoir'' and trains only
a \emph{linear} readout of its state~\cite{Maass2002,Jaeger2004}. RC \emph{requires}
dissipation: the echo-state / fading-memory property---that the influence of past
inputs decays---is supplied precisely by relaxation~\cite{Fujii2017,Nakajima2019}.
The molecular-spin and NMR communities have demonstrated quantum RC explicitly in
nuclear- and electron-spin ensembles~\cite{Fujii2017,Nakajima2019,Martinez2021},
part of a broad physical-reservoir-computing effort~\cite{Tanaka2019}.
In this system the fading memory is supplied not by electronic dephasing---removing
it entirely leaves the capacity and the memory kernel essentially unchanged---but by
the turnover itself: each cycle discards the electron pair and keeps only the
nuclei, which is a dissipative, non-unitary map. It is the openness of the chemical
cycle, not the decoherence of the electrons, that makes the system a reservoir.

Our contribution is to answer that question with the chemistry left in, and the
answer overturns three natural expectations. (i)~We drive the reservoir as a real
turnover does---nuclei retained, a new spin-correlated pair created each
cycle---rather than by resetting a single electron, and read it through the products
a downstream reaction could actually use. (ii)~We compare against classical
reservoirs matched by readout dimension, and find no quantum advantage, which we
report rather than argue around. (iii)~We show that the timescale objection usually
raised against this system rests on identifying the reservoir's clock with the
radical-pair lifetime, and replace it with a chemical criterion---reuse of the
nuclear register---that is sharper, falsifiable, and, unlike a comparison of two
timescales, actually requires the quantum treatment.

The physical correspondence we adopt, adapting the spin-RC framework of Fujii and
Nakajima~\cite{Fujii2017}, maps the reservoir-computing elements onto the radical
pair as follows: the long-$T_2$ memory nodes are the nuclear spins ($^{31}$P,
$^1$H); the nonlinear read/write head is the radical-pair electron system; the
read/write channel is the hyperfine coupling; and the fading memory (echo-state
property) is supplied by the electronic $T_2^e$ together with inter-cycle nuclear
dephasing. These are exactly the three layers of the hypothesis above: Layer~1
(nuclear-spin memory), Layer~2 (the radical-pair interface), and Layer~3 (the
classical readout/transduction that trains the linear map).

The paper is organized as follows. Section~\ref{sec:results} establishes the
reservoir's capacity and its convergence, localises that capacity in spin coherence,
compares it against matched classical baselines, and then asks what a chemical
readout recovers; it closes with the reservoir's clock and the nuclear-register
criterion. Section~\ref{sec:discussion} places the result among quantum-brain
proposals and states the limitations. The model, protocol, and capacity definitions
are collected in Appendix~\ref{app:methods}.

\section{Results}\label{sec:results}

\begin{figure}[htbp]
\includegraphics[width=\columnwidth]{fig_schematic.pdf}
\caption{\label{fig:schematic}
\textbf{The reservoir and the quantity that decides its reach.}
(a)~A cryptochrome-type separated radical pair ($J\approx0$, lifetime
$\tau\sim\mu$s) is the same molecule in two roles: long electronic coherence makes
it a weak-field magnetosensor, while the dissipation of an open cycle supplies the
fading memory a reservoir needs. (b)~The nuclei are the register: they are not
consumed by recombination and therefore persist from one turnover to the next,
whereas the electron pair is created afresh each cycle. The memory horizon is
consequently $\mathrm{MC}\times T_{\mathrm{turnover}}$, not
$\mathrm{MC}\times\tau$---and it reaches the neural band only if that register is
genuinely reused.}
\end{figure}

We drive the reservoir the way a real turnover does. Each cycle the nuclei are
retained and a \emph{new} spin-correlated electron pair is created whose singlet
character encodes the input,
$\rho_e(s)=s\,|S\rangle\langle S|+(1-s)P_T/3$; the pair then evolves under the
Liouville-space master equation with electronic dephasing and, where the readout
requires it, spin-selective recombination. This replaces the reset of a single
electron used in earlier treatments, which is not a chemical process and which---for
a separated pair---destroys the electron--electron correlation every cycle. All
capacities are out-of-sample (trained on the first half of a run, scored on the
held-out second half) and are means over $6$--$8$ input realisations. Model,
protocol and estimator are given in Appendix~\ref{app:methods}.

\subsection{Capacity of the reservoir}

With a twelve-observable readout the five-spin reservoir attains
\begin{equation*}
\mathrm{MC}=3.9,\qquad \mathrm{IPC}=5.6\pm0.2 ,
\end{equation*}
of which $1.7$ units are nonlinear. Unlike the single-electron-reset protocol, the
estimate is converged in the length of the driving sequence: repeating at
$L=600,1200,2400,4800$ gives $\mathrm{IPC}=5.28,5.60,5.87,5.84$, i.e.\ flat to
within the seed-to-seed scatter beyond $L\approx2400$. The memory kernel decays
over $\sim4$ cycles, the signature of the fading-memory property.

\subsection{The capacity is largely coherent}

To ask how much of this is quantum we add an all-spin pure-dephasing rate $\gamma$
in the $S_z$ product basis and take the incoherent limit at \emph{matched hopping
rate}: since the golden-rule transfer rate scales as $|H_{ij}|^2/\gamma$, holding
rate$\,\times\,\tau$ fixed requires $\tau\propto\gamma$. Doing so, the capacity
saturates at $\mathrm{IPC}=1.00$ for $\gamma=10^{3},5\times10^{3},2\times10^{4}$~MHz
--- exactly the memoryless floor, i.e.\ instantaneous read-back of the input with no
memory at all. Against the coherent value $5.42$ this is a coherence fraction
\begin{equation}
\frac{\mathrm{IPC}_{\gamma=0}-\mathrm{IPC}_{\rm incoh}}{\mathrm{IPC}_{\gamma=0}}
=0.82 .
\end{equation}
Taking the limit at fixed $\tau$ instead, as is sometimes done, measures Zeno
freezing rather than the classical limit and inflates this number; we avoid that.

\subsection{There is no quantum advantage}

A capacity is only meaningful against a classical reference, and the comparison must
be symmetric in what is read out. Driving tanh echo-state networks (ESNs) with the
same input and scoring them with the \emph{same} estimator:

\begin{center}
\begin{tabular}{lc}
\hline
reservoir (readout features) & IPC \\
\hline
quantum, 5 spins (12 observables) & $5.59\pm0.19$ \\
ESN, 5 nodes (5 linear features)  & $4.55\pm0.32$ \\
ESN, 5 nodes (12 features: $5$ linear $+7$ quadratic) & $9.41\pm0.86$ \\
ESN, 12 nodes (12 linear features) & $10.37\pm0.59$ \\
\hline
\end{tabular}
\end{center}

Matched by \emph{physical unit} the quantum reservoir is modestly ahead
($5.59$ vs $4.55$), but matched by \emph{readout dimension}---the resource that
actually sets the Dambre ceiling---a classical system of the same size is well ahead
($9.41$ vs $5.59$). We therefore make no claim of quantum advantage. Which
accounting is the fair one depends on a cost model for non-commuting measurements
versus multiplications that we do not have; we report the numbers and leave the
claim unmade.

\subsection{Biology can read the reservoir}

\begin{figure}[htbp]
\includegraphics[width=\columnwidth]{fig_readout.pdf}
\caption{\label{fig:readout}
\textbf{Readout is not the bottleneck.} Capacity recovered by chemically accessible
observables at the cryptochrome operating point ($B=50~\mu$T, $J=0$, $\tau=1~\mu$s),
with spin-selective recombination switched on so that the observable is the product
actually formed. The dotted line is the memoryless floor ($\mathrm{IPC}=1$): a
readout at that level reports the current input and nothing else. Ordinary reaction
kinetics clear it; adding the nuclear polarisation carried away by the product
(CIDNP) more than doubles the capacity again. No spectroscopy is involved.}
\end{figure}

Biology has no EPR spectrometer; it has products. We therefore restrict the readout
to observables a downstream reaction could plausibly use, with recombination
switched on so that the quantity read is the yield actually produced,
$Y_X=k_X\!\int_0^\tau\!\mathrm{Tr}[P_X\rho(t)]\,dt$ (Fig.~\ref{fig:readout}):

\begin{center}
\begin{tabular}{lcc}
\hline
readout route & channels & IPC \\
\hline
end-point singlet yield & 1 & $1.02$ \\
both exit channels, end point & 2 & $1.37$ \\
three micro-environments (different $\tau$) & 3 & $2.56$ \\
time-resolved kinetics of one product pool & 5 & $2.01$ \\
\quad + accumulated downstream pool & 10 & $2.71$ \\
\quad + CIDNP polarisation on the product & 8 & $\mathbf{4.65}$ \\
\hline
\end{tabular}
\end{center}

Two things follow. First, the readout is not the limiting resource: time-resolved
product kinetics alone clear the memoryless floor by a factor of two, and the
nuclear polarisation the product carries away raises the capacity to $4.65$ without
any spectroscopic access whatsoever. Second, a single scalar end-point yield gives
$\mathrm{IPC}=1.02$, which saturates its own one-observable bound but consists
almost entirely of the zero-delay term: it reports the present input and remembers
nothing. Claims about reservoir behaviour should not be based on it.

The mechanism behind the CIDNP column is worth stating, because it also explains a
negative result. The reduced state of \emph{either} electron of a newly born pair is
maximally mixed for every input, since both $|S\rangle\langle S|$ and $P_T/3$ have
zero one-electron polarisation; the input lives entirely in the two-electron
correlation. With $J=0$ the two radicals do not interact, so nothing can transfer
that correlation to a nucleus---and indeed, with recombination switched off, the
cryptochrome-point capacity is $1.01$, the memoryless floor. It is
\emph{spin-selective recombination} that converts the singlet--triplet correlation
into a population difference and so writes the input into the nuclear register. In a
separated radical pair, recombination is not merely how the reservoir is read; it is
how it is written.

\subsection{The clock is the turnover, not the pair lifetime}

\begin{figure}[htbp]
\includegraphics[width=\columnwidth]{fig_criterion.pdf}
\caption{\label{fig:criterion}
\textbf{(a)} The memory horizon $\mathrm{MC}\times T_{\rm turnover}$ against the
turnover interval. Because the nuclear register survives recombination, inserting a
physiological pause leaves the memory capacity essentially unchanged and moves the
horizon into the neural band (shaded). The dotted line is where the horizon would
sit if the clock were the radical-pair lifetime.
\textbf{(b)} The criterion. As the nuclear register is depolarised between
turnovers with probability $q$, the capacity falls to the memoryless floor: at
$q=1$---a fresh molecule each cycle---$\mathrm{MC}=1.00$ and there is no memory.}
\end{figure}

The horizon over which a reservoir remembers is $\mathrm{MC}$ multiplied by its
cycle time. It is natural, and it is what the submitted version of this work did, to
identify that cycle time with the radical-pair lifetime $\tau\sim1~\mu$s and
conclude that a $3~\mu$s horizon puts the system $10^3$--$10^5$ below neural
timescales. That identification is wrong, and its own model says so. The carrier of
the memory is the nuclear register, and nuclei are not consumed by recombination:
they leave in the diamagnetic product, where nuclear relaxation times are of order
$0.1$--$1$~s. The cycle time of the reservoir is therefore the interval between
turnovers.

Inserting a pause $T_{\rm d}$ between turnovers and relaxing the register with
probability $q=1-\exp(-T_{\rm d}/T_{1n})$ (with $T_{1n}=1$~s) gives
Fig.~\ref{fig:criterion}(a):

\begin{center}
\begin{tabular}{lccc}
\hline
$T_{\rm d}$ & $q$ & MC (8 ch) & horizon \\
\hline
$1~\mu$s & $10^{-6}$ & $2.818$ & $5.6~\mu$s \\
$1$~ms & $10^{-3}$ & $2.818$ & $2.8$~ms \\
$10$~ms & $10^{-2}$ & $2.817$ & $28$~ms \\
$100$~ms & $9.5\times10^{-2}$ & $2.800$ & $280$~ms \\
\hline
\end{tabular}
\end{center}

The memory capacity changes by $0.2\%$ across five orders of magnitude in the
turnover interval, so the horizon simply tracks it, and a perfectly ordinary
enzymatic or photocycle turnover of $1$--$100$~ms places it squarely in the neural
band. The timescale argument does not survive; the same spin physics runs on a slow
clock.

\subsection{The criterion: reuse of the nuclear register}

What replaces it is a chemical requirement, and a sharper one. Everything above
assumes that the register the reservoir writes into is the \emph{same} register on
the next turnover. Relaxing that assumption by depolarising the nuclei with
probability $q$ between cycles gives Fig.~\ref{fig:criterion}(b):
memory is essentially intact for $q\lesssim0.5$ (that is, while the turnover
interval is shorter than the nuclear relaxation time), degrades through
$q\approx0.7$--$0.9$, and at $q=1$---a fresh molecule every cycle---collapses to
$\mathrm{MC}=1.00$, the memoryless floor, for every readout we tested.

This is the constraint that actually decides whether a radical pair can serve as a
biological reservoir:

\begin{quote}
\emph{The same nuclear register must be re-used across successive turnovers, and
successive turnovers must be driven by distinct input samples. Given that, the
memory horizon is $\mathrm{MC}\times T_{\rm turnover}$, bounded above by
$\mathrm{MC}\times T_{1n}\sim1$~s; without it, the system retains nothing.}
\end{quote}

Unlike a comparison of two timescales---which requires no quantum mechanics at
all---this criterion is a statement about a spin register and can only be evaluated
by the kind of analysis performed here. It is also directly falsifiable by chemistry:
it fails if the radical pair is regenerated on a fresh molecule each turnover, or if
the relevant nuclear relaxation in the diamagnetic intermediate is faster than the
turnover rate.

\section{Discussion}\label{sec:discussion}

\paragraph*{What changed, and why it matters.}
Three claims that are natural to make about this system---that a scalar chemical
readout is too impoverished to be useful, that a spin reservoir of a few spins beats
a classical system of the same size, and that microsecond radical-pair lifetimes put
biological computation out of reach---all fail when the model is driven as a real
turnover and compared symmetrically. The readout is not the bottleneck; the quantum
advantage is not there; and the timescale gap is an artifact of identifying the
reservoir clock with the pair lifetime rather than the turnover interval. What
remains is a single, sharp, chemical criterion.

\paragraph*{Relation to prior work.}
Quantum reservoir computing in spin ensembles is established, including the use of
dissipation as a resource and the exponential-state-space
argument.\cite{Fujii2017,Nakajima2019,Martinez2021,Mujal2021,Tanaka2019} We do not
claim these as new, and---having found no advantage under a symmetric
comparison---we do not add to them. Our contribution is to take a chemically
realised radical pair, drive and read it as chemistry actually would, and convert a
long-standing qualitative conjecture\cite{Fisher2015,Kerskens2022,Tegmark2000} into a
criterion that experiment can test. The result also sharpens why CIDNP matters here:
in a separated pair it is not an optional extra channel but the mechanism by which
information reaches the nuclei at all.

\paragraph*{Limitations.}
The model is a five-spin reduction with isotropic hyperfine couplings; the full
anisotropic tensor set of the cryptochrome pair\cite{Lee2014,Schleicher2021} is not
treated, and $^{14}$N is modelled as spin-$1/2$. The capacity estimator uses
degree-$2$ targets, so higher-order nonlinear capacity is not counted. Whether a
\emph{neuronal} cryptochrome forms a separated radical pair at all is
unestablished---CRY1/CRY2 are primarily circadian transcriptional
regulators\cite{Partch2014}---so the biological reading of our criterion is
conditional. The spin-dynamics results do not depend on that condition.

\paragraph*{What would settle it.}
The criterion is addressable in vitro before it is addressable in a neuron.
Covalently linked donor--bridge--acceptor radical pairs have synthetically tunable
exchange coupling and lifetime and are read by transient EPR and optical/CIDNP
spectroscopy.\cite{Wasielewski2006,Wasielewski2020,Hore2016} Repeatedly cycling such
a system and asking whether the product yield of turnover $k$ depends on the input
at turnover $k-1$, $k-2$, \ldots{} measures the memory horizon directly, and varying
the interval between cycles tests the register-reuse criterion where it is sharpest.

\section{Conclusion}

Driven and read as a chemical turnover rather than as an abstract quantum channel, a
cryptochrome-type radical pair is a reservoir whose output ordinary chemistry can
recover: time-resolved product kinetics give $\mathrm{IPC}=2.0$ and the nuclear
polarisation carried off by the product raises this to $4.7$, with no spectroscopic
access required. Eighty-two per cent of the capacity is coherent in origin, and yet
there is no quantum advantage---a classical network with the same number of readout
features matches or exceeds it. The obstacle usually invoked, a microsecond memory
horizon, does not survive scrutiny: the register is nuclear and survives
recombination, so the clock is the turnover interval and a physiological
$1$--$100$~ms turnover places the horizon at $3$--$280$~ms. What decides the question
instead is whether the same nuclear register is reused from one turnover to the
next; if it is not, the capacity falls to a memoryless $\mathrm{MC}=1$. We offer that
criterion, rather than a timescale comparison, as the thing to test.

\begin{acknowledgments}
We thank the reviewers of the companion sensing study, whose insistence on
decoherence-explicit treatment shaped the methodology used here. All numerical
results are reproducible from the open-source \texttt{qbscreen} package.
\end{acknowledgments}

\section*{Author contributions}
H.W. conceived the study, developed the master-equation and reservoir code, ran the
simulations, and wrote the manuscript. T.T. contributed to the spin-dynamics
formulation and the analysis of the nuclear-register criterion. Both authors
reviewed the results and approved the manuscript.

\appendix

\section{Model and methods}
\label{app:methods}

\subsection{Reservoir Hamiltonian}
We model the reservoir as $N=5$ spins---two radical-pair electrons
$(e_1,e_2)$ and three nuclei ($^{31}$P and two $^1$H)---with Hamiltonian
(in MHz)
\begin{align}
H &= \omega_e\,(S_z^{e_1}+S_z^{e_2})
   + \sum_{\mu=x,y,z}\Big[
     A_P\,S_\mu^{e_1}I_\mu^{P}
   + A_{H_1}S_\mu^{e_1}I_\mu^{H_1} \nonumber\\
&\quad
   + A_{e_2}S_\mu^{e_2}I_\mu^{H_2}
   + J\,S_\mu^{e_1}S_\mu^{e_2}\Big],
\label{eq:H}
\end{align}
where $\omega_e=g_e\mu_B B/h$ is the electron Zeeman frequency. Each radical carries
its own nuclei---they are not shared, which is what a spatially separated pair
means---and the exchange $J$ links the two electrons. At the engineered operating
point we use $A_{e_1a}=80$, $A_{e_1b}=40$, $A_{e_2a}=15$~MHz, $J=2$~MHz and
$B=1$~mT: representative molecular-spin-RC values chosen to give effective hyperfine
writing, not a fit to a specific protein. For the cryptochrome point we set $J=0$,
$B=50~\mu$T and $\tau=1~\mu$s, with hyperfine magnitudes representative of the
dominant FAD$^{\bullet-}$ and TrpH$^{\bullet+}$ couplings~\cite{Lee2014,Schleicher2021}.
The conclusions are not fine-tuned to these numbers: the capacity varies by less
than $\sim20\%$ over an order-of-magnitude scan of $T_2^e$, and the register
criterion of Sec.~\ref{sec:results} is independent of the hyperfine assignment.

\subsection{Master equation and propagator}
Each reservoir update is the completely positive, trace-preserving map
$\rho\mapsto e^{\mathcal{L}\tau}\rho$ generated by the Haberkorn master
equation~\cite{Haberkorn1976} (whose spin-selective recombination operator we adopt
in its standard form~\cite{Jones2010}) with electronic pure dephasing,
\begin{equation}
\frac{d\rho}{dt}
 = -i[H,\rho]
 - \tfrac12\{k_S P_S + k_T P_T,\rho\}
 + \sum_{i}\frac{1}{T_2^e}\!\left(S_z^i\rho S_z^i-\tfrac14\rho\right),
\label{eq:me}
\end{equation}
solved exactly in Liouville space. The third term---the electronic decoherence
that the companion sensing analysis showed to be decisive---is here the source of
fading memory. We verified that, in the limit $1/T_2^e\to0$ and $k_S=k_T$, the
Liouville-space solver reproduces the closed-form coherent Green's-function singlet
yield to machine precision (Supplemental Material). Within an update step no
recombination acts ($k_S=k_T=0$); the readout step uses the spin-selective term.

\subsection{Reservoir protocol}
We use the standard temporal-QRC protocol~\cite{Fujii2017}. At step $k$ the input
electron is reset to the input-encoded pure state
$|\psi(s_k)\rangle=\sqrt{1-s_k}\,|0\rangle+\sqrt{s_k}\,|1\rangle$ while the
reservoir (the remaining electron plus nuclei) is retained; the full state evolves
for a fixed lifetime $\tau$ under $e^{\mathcal{L}\tau}$ ($\tau=50$~ns,
$T_2^e=50$~ns at the operating point); a set of observables is read out, and a
linear (ridge) readout is trained on them. Inputs $s_k$ are i.i.d.\ uniform on
$[0,1]$; the first $\sim100$ steps are discarded as washout.

This protocol has a natural chemical reading that addresses its physical
realizability. The ``reset'' of the input electron is not an applied operation but
the \emph{birth of a new radical pair}: each catalytic or photochemical cycle
creates a fresh spin-correlated electron pair whose initial state encodes the
instantaneous input (e.g.\ a substrate-flux- or light-dependent singlet/triplet
bias), while the surrounding nuclei---which are not consumed by recombination---carry
spin polarisation across cycles. The reservoir is thus a \emph{persistent nuclear
bath repeatedly interrogated by transient electron pairs}, which is precisely the
radical-pair situation, not an artificial gate sequence. The washout discards
initial transients, so the reported capacities are independent of the initial spin
state---a direct consequence of the fading-memory (echo-state) property and the
reason no thermal-population fine-tuning is needed.

\subsection{Capacity measures}
We quantify the reservoir with two information-theoretic benchmarks. The linear
\emph{memory capacity}~\cite{Jaeger2002}
$\mathrm{MC}=\sum_{d\ge0} r^2(s_{k-d})$ sums the squared multiple correlation for
reconstructing delayed inputs. The \emph{information processing
capacity}~\cite{Dambre2012} adds nonlinear functions of the input history,
$\mathrm{IPC}=\mathrm{MC}+\sum_{\text{nonlinear}}C$, evaluated here through
degree-2 Legendre targets (diagonal $P_2(s_{k-d})$ and cross
$s_{k-d_1}s_{k-d_2}$ terms). A fundamental theorem~\cite{Dambre2012} bounds the
total capacity by the number of linearly independent readout observables,
$\mathrm{IPC}\le n_\mathrm{obs}$; we use this bound throughout as a correctness
anchor and as the key to interpreting the ensemble results.

All results are reproducible from the open \texttt{qbscreen} package
(\texttt{reservoir.py}, \texttt{ensemble.py}, \texttt{quantum\_vs\_classical.py};
28/28 tests pass).


\section{Data and code availability}
All analyses are implemented in the open \texttt{qbscreen} package,
\texttt{https://github.com/deeptell-inc/new\_Q\_brain}. The corrected
correlated-pair-birth injection is in \texttt{qbscreen/corrected\_injection.py}; the
chemically accessible readout routes, the turnover-clock scan and the
register-reuse criterion are in \texttt{qbscreen/readout\_routes.py}; the
protocol-consistent capacity, coherence fraction and classical baselines are in
\texttt{qbscreen/final\_numbers.py}; the master-equation solver is in
\texttt{qbscreen/master\_equation.py}. Numerical outputs are in
\texttt{simulation\_results/}. The capacity bound and the agreement of the solver
with the closed-form coherent limit are checked in \texttt{qbscreen/tests/}
(28 tests, all passing).  
           
\begin{thebibliography}{99}  
          
\bibitem{Steiner1989} U.~E. Steiner and T.~Ulrich,
\textit{Magnetic field effects in chemical kinetics and related phenomena},
Chem. Rev. \textbf{89}, 51 (1989).
 
\bibitem{Hore2016} P.~J. Hore and H.~Mouritsen,
\textit{The radical-pair mechanism of magnetoreception},
Annu. Rev. Biophys. \textbf{45}, 299 (2016).

\bibitem{Ritz2000} T.~Ritz, S.~Adem, and K.~Schulten,
\textit{A model for photoreceptor-based magnetoreception in birds},
Biophys. J. \textbf{78}, 707 (2000).

\bibitem{Maeda2012} K.~Maeda \textit{et al.},
\textit{Magnetically sensitive light-induced reactions in cryptochrome are
consistent with its proposed role as a magnetoreceptor},
Proc. Natl. Acad. Sci. USA \textbf{109}, 4774 (2012).

\bibitem{Xu2021} J.~Xu \textit{et al.},
\textit{Magnetic sensitivity of cryptochrome 4 from a migratory songbird},
Nature \textbf{594}, 535 (2021).

\bibitem{Fisher2015} M.~P.~A. Fisher,
\textit{Quantum cognition: The possibility of processing with nuclear spins in the
brain}, Ann. Phys. \textbf{362}, 593 (2015).

\bibitem{Kerskens2022} C.~M. Kerskens and D.~López Pérez,
\textit{Experimental indications of non-classical brain functions},
J. Phys. Commun. \textbf{6}, 105001 (2022).

\bibitem{Buzsaki2004} G.~Buzsáki and A.~Draguhn,
\textit{Neuronal oscillations in cortical networks},
Science \textbf{304}, 1926 (2004).

\bibitem{Tegmark2000} M.~Tegmark,
\textit{Importance of quantum decoherence in brain processes},
Phys. Rev. E \textbf{61}, 4194 (2000).

\bibitem{Lambert2013} N.~Lambert, Y.-N. Chen, Y.-C. Cheng, C.-M. Li, G.-Y. Chen,
and F.~Nori, \textit{Quantum biology}, Nat. Phys. \textbf{9}, 10 (2013).

\bibitem{Cao2020} J.~Cao \textit{et al.},
\textit{Quantum biology revisited}, Sci. Adv. \textbf{6}, eaaz4888 (2020).

\bibitem{Hiscock2016} H.~G. Hiscock \textit{et al.},
\textit{The quantum needle of the avian magnetic compass},
Proc. Natl. Acad. Sci. USA \textbf{113}, 4634 (2016).

\bibitem{Tanaka2019} G.~Tanaka \textit{et al.},
\textit{Recent advances in physical reservoir computing: A review},
Neural Netw. \textbf{115}, 100 (2019).

\bibitem{Wasielewski2020} M.~R. Wasielewski \textit{et al.},
\textit{Exploiting chemistry and molecular systems for quantum information science},
Nat. Rev. Chem. \textbf{4}, 490 (2020).

\bibitem{Partch2014} C.~L. Partch, C.~B. Green, and J.~S. Takahashi,
\textit{Molecular architecture of the mammalian circadian clock},
Trends Cell Biol. \textbf{24}, 90 (2014).

\bibitem{Timmel1998} C.~R. Timmel, U.~Till, B.~Brocklehurst, K.~A. McLauchlan,
and P.~J. Hore, \textit{Effects of weak magnetic fields on free radical
recombination reactions}, Mol. Phys. \textbf{95}, 71 (1998).

\bibitem{Jones2010} J.~A. Jones and P.~J. Hore,
\textit{Spin-selective reactions of radical pairs act as quantum measurements},
Chem. Phys. Lett. \textbf{488}, 90 (2010).

\bibitem{HoreBroadhurst1993} P.~J. Hore and R.~W. Broadhurst,
\textit{Photo-CIDNP of biopolymers},
Prog. Nucl. Magn. Reson. Spectrosc. \textbf{25}, 345 (1993).

\bibitem{Wakaura2025} H.~Wakaura,
\textit{Systematic screening of brain neurotransmitter enzymes for radical-pair
mechanism competence identifies candidates for magnetic field sensitivity},
Research Square preprint, doi:10.21203/rs.3.rs-9278975/v1 (2025).

\bibitem{Wakaura2026} H.~Wakaura,
\textit{Motional narrowing rescues entanglement in biological radical-pair systems:
quantum error correction and reservoir computing beyond the breaking threshold},
Research Square preprint, doi:10.21203/rs.3.rs-9493052/v1 (2026).

\bibitem{Lee2014} A.~A. Lee, J.~C.~S. Lau, H.~J. Hogben, T.~Biskup, D.~R. Kattnig,
and P.~J. Hore,
\textit{Alternative radical pairs for cryptochrome-based magnetoreception},
J. R. Soc. Interface \textbf{11}, 20131063 (2014).

\bibitem{Schleicher2021} E.~Schleicher \textit{et al.},
\textit{Selective $^{13}$C labelling reveals the electronic structure of
flavocoenzyme radicals}, Sci. Rep. \textbf{11}, 18395 (2021).

\bibitem{Efimova2008} O.~Efimova and P.~J. Hore,
\textit{Role of exchange and dipolar interactions in the radical pair model of the
avian magnetic compass}, Biophys. J. \textbf{94}, 1565 (2008).

\bibitem{Companion2026} H.~Wakaura and T.~Tanimae,
\textit{Electronic decoherence and exchange coupling preclude weak-field
magnetosensitivity of neurotransmitter-metabolising flavoenzymes}
(companion manuscript, 2026).

\bibitem{Maass2002} W.~Maass, T.~Natschläger, and H.~Markram,
\textit{Real-time computing without stable states: A new framework for neural
computation based on perturbations}, Neural Comput. \textbf{14}, 2531 (2002).

\bibitem{Jaeger2004} H.~Jaeger and H.~Haas,
\textit{Harnessing nonlinearity: Predicting chaotic systems and saving energy in
wireless communication}, Science \textbf{304}, 78 (2004).

\bibitem{Fujii2017} K.~Fujii and K.~Nakajima,
\textit{Harnessing disordered-ensemble quantum dynamics for machine learning},
Phys. Rev. Applied \textbf{8}, 024030 (2017).

\bibitem{Nakajima2019} K.~Nakajima \textit{et al.},
\textit{Boosting computational power through spatial multiplexing in quantum
reservoir computing}, Phys. Rev. Applied \textbf{11}, 034021 (2019).

\bibitem{Martinez2021} R.~Martínez-Peña \textit{et al.},
\textit{Dynamical phase transitions in quantum reservoir computing},
Phys. Rev. Lett. \textbf{127}, 100502 (2021).

\bibitem{Mujal2021} P.~Mujal, R.~Martínez-Peña, J.~Nokkala, J.~García-Beni,
G.~L. Giorgi, M.~C. Soriano, and R.~Zambrini,
\textit{Opportunities in quantum reservoir computing and extreme learning
machines}, Adv. Quantum Technol. \textbf{4}, 2100027 (2021).

\bibitem{Haberkorn1976} R.~Haberkorn,
\textit{Density matrix description of spin-selective radical pair reactions},
Mol. Phys. \textbf{32}, 1491 (1976).

\bibitem{Jaeger2002} H.~Jaeger,
\textit{Short term memory in echo state networks},
GMD Report 152, German National Research Center for Information Technology (2002).

\bibitem{Dambre2012} J.~Dambre, D.~Verstraeten, B.~Schrauwen, and S.~Massar,
\textit{Information processing capacity of dynamical systems},
Sci. Rep. \textbf{2}, 514 (2012).

\bibitem{Wasielewski2006} M.~R. Wasielewski,
\textit{Energy, charge, and spin transport in molecules and self-assembled
nanostructures inspired by photosynthesis},
J. Org. Chem. \textbf{71}, 5051 (2006).

\end{thebibliography}

\end{document}
